% Unit 072 terjemahan Bahasa Indonesia dari chapter6.tex baris 1--320.
% Sumber: Methods of Algebra, Volume 2, commit
% 9a5803ff77dd3257484cb177f851a73770a59dd3 (CC BY 4.0).
% stable-unit-id: o014.aljabr2.chapter6.overview-g-modules-resolutions
% Cakupan: pembukaan Bab 6 dan Bagian 6.1 lengkap tentang modul-G serta
% resolusi bebas standar dan ternormalisasi; berhenti tepat sebelum Bagian 6.2
% pada baris 321.

% segment-id: o014.li2.chapter6.overview-g-modules-resolutions.h001
\chapter{Homologi dan Kohomologi Grup}\label{sec:group-coh}
\hypertarget{o014.aljabr2.chapter6.overview-g-modules-resolutions}{ }

% segment-id: o014.li2.chapter6.overview-g-modules-resolutions.p001
Misalkan $G$ suatu grup dan $\Bbbk$ suatu gelanggang komutatif. Yang dimaksud
dengan modul-$G$ ialah modul-$\Bbbk$ yang dilengkapi aksi kiri $G$ yang linear;
aksinya ditulis sebagai perkalian kiri. Dalam banyak situasi, $\Bbbk$ biasanya
diambil sama dengan $\Z$. Semua modul-$G$ membentuk kategori abelian
$G\dcate{Mod}$, yang isomorfik dengan $\Bbbk[G]\dcate{Mod}$, dengan
$\Bbbk[G]$ aljabar grup dari $G$. Untuk modul-$G$ $M$, definisikan
modul-$\Bbbk$
% segment-id: o014.li2.chapter6.overview-g-modules-resolutions.q001
\begin{align*}
	\text{invarian} \quad M^G & := \left\{x \in M: \forall g \in G, \; gx=x \right\}, \\
	\text{koinvarian} \quad M_G & := M / \lrangle{gx-x: g \in G, x \in M}.
\end{align*}

% segment-id: o014.li2.chapter6.overview-g-modules-resolutions.p002
Dari sudut pandang aljabar, kohomologi $\Hm^n(G,M)$ dan homologi
$\Hm_n(G,M)$ dari modul-$G$ $M$ masing-masing merupakan nilai pada $M$ dari
funktor turunan kanan $(\cdot)^G$ dan funktor turunan kiri $(\cdot)_G$, dengan
$n\in\Z_{\geq0}$.

% segment-id: o014.li2.chapter6.overview-g-modules-resolutions.p003
Homologi dan kohomologi grup memiliki sejarah yang cukup panjang. Sekitar tahun
1935, W.\ Hurewicz mengkaji ruang topologis tersambung lintasan $E$ yang grup
homotopi tingkat tingginya memenuhi
$n>1\implies\pi_n(E)=0$. Ia membuktikan bahwa grup homologi dan kohomologinya
sepenuhnya ditentukan oleh grup fundamental $\pi_1(E)$; ini mungkin merupakan
rumusan kohomologi grup paling awal yang tercatat. Seiring kemajuan pesat
aljabar homologis pada abad ke-20, kohomologi dan homologi grup segera
memperoleh rumusan aljabar murni dan digunakan untuk menangani persoalan dalam
aljabar itu sendiri. Dewasa ini, keduanya telah menjadi perangkat dasar dalam
teori bilangan, topologi, dan geometri.

% segment-id: o014.li2.chapter6.overview-g-modules-resolutions.p004
Di sisi lain, $\Hm^n(G,M)$ dan $\Hm_n(G,M)$ masing-masing dapat diidentifikasi
dengan $\Ext^n(\Bbbk,M)$ dan $\Tor_n(\Bbbk,M)$ dalam
$\Bbbk[G]\dcate{Mod}$, dengan $\Bbbk$ dipandang sebagai modul-$G$ trivial.
Dengan menghitung $\Ext$ dan $\Tor$ melalui suatu resolusi projektif dari
$\Bbbk$, diperoleh deskripsi lain bagi $\Hm^n(G,M)$ dan $\Hm_n(G,M)$. Bahkan,
kita akan memberikan resolusi kanonik
$\mathsf{L}\to\Bbbk$ bagi $\Bbbk$ sebagai modul kiri (atau kanan)
$\Bbbk[G]$. Dari resolusi tersebut diperoleh kompleks kokrantai standar
$(C^n(G,M))_n$ (atau kompleks rantai $(C_n(G,M))_n$) yang menyajikan
$\Hm^n(G,M)$ (atau $\Hm_n(G,M)$); lihat
Proposisi~\ref{prop:groupcoh-cochain} dan \ref{prop:grouph-chain}. Selanjutnya,
\S\S\ref{sec:homology-computation}---\ref{sec:bar-resolution} akan menafsirkan
resolusi kanonik ini dari sudut pandang teori simpleks.

% segment-id: o014.li2.chapter6.overview-g-modules-resolutions.p005
Hal-hal di atas merupakan isi utama
\S\S\ref{sec:G-mod}---\ref{sec:group-coh-sub}. Bab ini memakai dua sudut
pandang sekaligus: teori klasik funktor turunan (khususnya funktor-$\delta$
universal beserta karakterisasinya) dan kategori turunan. Kelebihan teori
klasik ialah nuansa operasionalnya, sedangkan kategori turunan menyediakan
informasi yang lebih halus daripada tiap-tiap grup kohomologi atau homologi dan
sering kali juga menawarkan gagasan yang lebih jernih. Kompleks kokrantai atau
rantai standar memberikan informasi yang paling terperinci dan kanonik, sebab
konstruksi itu bekerja pada tingkat kompleks.

% segment-id: o014.li2.chapter6.overview-g-modules-resolutions.p006
Dalam \S\ref{sec:grp-low-degree}, kita akan memberikan penafsiran konkret bagi
$\Hm^1$ dan $\Hm^2$: keduanya berturut-turut bersesuaian dengan homomorfisme
silang dan kelas ekuivalensi ekstensi $G$ oleh grup abelian. Makna terperinci
ekstensi grup dapat dilihat pada Definisi~\ref{def:group-ext}. Bahkan, jika
semua ekstensi grup dari $G$ oleh $M$ disusun menjadi 2-kategori
$\cate{Ext}(G,M)$, pemenggalan
$\tau^{\leq2}\overline{C}(G,M)$ dari kompleks standar ternormalisasi tepat
merupakan suatu perwujudan linear dari $\cate{Ext}(G,M)$; inilah isi
Catatan~\ref{rem:Ext-incarnation}.

% segment-id: o014.li2.chapter6.overview-g-modules-resolutions.p007
Untuk sembarang homomorfisme grup $\varphi:H\to G$ dan modul-$G$ $M$, dapat
didefinisikan tarikan balik $\varphi^*M$ sebagai modul-$H$. Berdasarkan
konstruksi umum dalam teori modul, adjoin kanan dan adjoin kiri dari
$\varphi^*$ dapat dituliskan dengan mudah. Jika $\varphi$ merupakan inklusi
suatu subgrup, $\varphi^*$ disebut funktor restriksi modul-$G$ dan juga ditulis
$\Res^G_H$, sedangkan adjoin kanan (atau kirinya) ditulis $\Ind^G_H$ (atau
$\iInd^G_H$); keduanya disebut funktor induksi. Sifat-sifatnya akan menjadi
pokok bahasan \S\ref{sec:induced-module}. Secara khusus, kita akan membuktikan
apa yang disebut Lema Shapiro (Teorema~\ref{prop:Shapiro}):
% segment-id: o014.li2.chapter6.overview-g-modules-resolutions.q002
\[ \Hm^n\left(G, \Ind^G_H(N)\right) \simeq \Hm^n(H, N), \quad \Hm_n\left(G, \iInd^G_H(N)\right) \simeq \Hm_n(H, N). \]

% segment-id: o014.li2.chapter6.overview-g-modules-resolutions.p008
Melanjutkan pembahasan umum tentang $\varphi^*$,
\S\ref{sec:change-of-group} akan mendefinisikan keluarga homomorfisme kanonik
% segment-id: o014.li2.chapter6.overview-g-modules-resolutions.q003
\[ \Hm^n(G, M) \to \Hm^n(H, \varphi^* M), \quad \Hm_n(H, \varphi^* M) \to \Hm_n(G, M). \]
Homomorfisme ini dapat dipadukan dengan funktorialitas kohomologi atau homologi,
sehingga setiap homomorfisme ekuivarian $f:M\to N$ (atau $f:N\to M$)
menginduksi homomorfisme kanonik yang bersesuaian, yaitu $\Hm^n(f)$ (atau
$\Hm_n(f)$), dengan $M$ modul-$G$ dan $N$ modul-$H$; lihat
Definisi~\ref{def:change-of-group-equiv}. Untuk kasus khusus $n=2$,
\S\ref{sec:group-ext-revisited} akan menafsirkan homomorfisme kanonik tersebut
dari sudut pandang ekstensi grup.

% segment-id: o014.li2.chapter6.overview-g-modules-resolutions.p009
Untuk memperoleh gambaran konkret, \S\ref{sec:coh-cyclic-free} membahas contoh
grup siklik hingga dan grup bebas. Argumen yang digunakan memadukan teknik
teori grup dan kohomologi secara serasi. Dimensi kohomologis dan
dimensi homologis grup (Definisi~\ref{def:group-dim}) juga akan muncul pada
saat yang tepat.

% segment-id: o014.li2.chapter6.overview-g-modules-resolutions.p010
Nuansa konkret teori grup berlanjut dalam \S\ref{sec:group-finite-index}. Untuk
kasus $(G:H)$ hingga, kita akan menyatakan sifat-sifat lebih lanjut dari funktor
restriksi dan induksi; bahkan, dalam keadaan ini
$\Ind^G_H=\iInd^G_H$. Definisi--Proposisi~\ref{def:cor} akan memberikan
keluarga homomorfisme kanonik
% segment-id: o014.li2.chapter6.overview-g-modules-resolutions.q004
\[ \mathrm{cor}^n: \Hm^n(H, M) \to \Hm^n(G, M), \quad \mathrm{cor}_n: \Hm_n(G, M) \to \Hm_n(H, M), \]
dengan $M$ suatu modul-$G$ dan $\Res^G_HM$ juga disingkat menjadi $M$ untuk
menghemat notasi. Homomorfisme ini disebut korestriksi dan arahnya berlawanan
dengan homomorfisme kanonik yang diinduksi oleh $H\hookrightarrow G$ (yaitu
restriksi, yang ditulis $\mathrm{res}^n$ dan $\mathrm{res}_n$). Keduanya
memenuhi identitas dasar
$\mathrm{cor}^n\mathrm{res}^n=(G:H)$ dan
$\mathrm{res}_n\mathrm{cor}_n=(G:H)$
(Proposisi~\ref{prop:res-cor}), sedangkan kasus khusus
$\mathrm{cor}_1:\Hm_1(G,\Z)\to\Hm_1(H,\Z)$ menjadi homomorfisme transfer
klasik dalam teori grup,
$\mathrm{Ver}_{G|H}:G_{\mathrm{ab}}\to H_{\mathrm{ab}}$.

% segment-id: o014.li2.chapter6.overview-g-modules-resolutions.p011
Misalkan $H$ subgrup normal dari $G$ dan $M$ modul-$G$. Dalam keadaan ini,
$\Hm^q(H,M)$ dan $\Hm_q(H,M)$ menjadi modul-$(G/H)$. Barisan spektral
Lyndon--Hochschild--Serre yang dibahas dalam \S\ref{sec:LHS-SS},
% segment-id: o014.li2.chapter6.overview-g-modules-resolutions.q005
\begin{align*}
	E_2^{p, q} = \Hm^p(G/H, \Hm^q(H, M)) & \Rightarrow \Hm^{p+q}(G, M), \\
	E^2_{p, q} = \Hm_p(G/H, \Hm_q(H, M)) & \Rightarrow \Hm_{p+q}(G, M),
\end{align*}
merupakan perangkat tingkat lanjut yang mutlak diperlukan dalam teori
kohomologi dan homologi grup; lihat Teorema~\ref{prop:LHS-SS}. Informasi dalam
barisan eksak suku-suku berderajat rendahnya sangat berguna. Barisan spektral
Lyndon--Hochschild--Serre sekaligus merupakan penerapan barisan spektral
Grothendieck dan berasal dari suatu filtrasi kanonik pada kompleks standar.
Deskripsi terakhir memberikan informasi yang lebih halus, misalnya struktur
multiplikatif yang disebutkan dalam Catatan~\ref{rem:LHS-SS-mult}.

% segment-id: o014.li2.chapter6.overview-g-modules-resolutions.p012
Lebih tepatnya, operasi multiplikatif pada kohomologi grup ialah hasil kali
cawan yang berasal dari topologi. Dalam \S\ref{sec:cup-product}, kita akan
mendekatinya dari dua sudut pandang. Pertama, hasil kali cawan dipahami sebagai
pencerminan alami operasi hasil kali tensor modul-$G$; bahasa kategori turunan
membantu dalam hal ini. Kedua, rumusnya dituliskan langsung pada tingkat
kompleks standar (Definisi--Proposisi~\ref{def:group-cup-3}). Deskripsi kedua
tentu merupakan perwujudan eksplisit dari deskripsi pertama. Kunci penurunannya
ialah mengangkat $\Bbbk\rightiso\Bbbk\otimes\Bbbk$ secara tepat menjadi
``pembenaman diagonal''
$\Delta:\mathsf{L}\to\mathsf{L}\otimes\mathsf{L}$; beberapa tanda dalam
argumen ini perlu ditangani dengan hati-hati.

% segment-id: o014.li2.chapter6.overview-g-modules-resolutions.p013
Jika $G$ hingga, kohomologi Tate $\TaHm^n(G,M)$ yang diperkenalkan dalam
\S\ref{sec:Tate-coh} memadukan kohomologi dan homologi. Konstruksinya bergantung
pada homomorfisme kanonik yang khas bagi grup hingga,
% segment-id: o014.li2.chapter6.overview-g-modules-resolutions.q006
\[ \nu: M_G \to M^G, \quad (\text{citra dari }x \in M) \mapsto \sum_{g \in G} gx. \]
Untuk $n\geq1$ (atau $n\leq-2$), kohomologi Tate sama dengan $\Hm^n(G,M)$
(atau $\Hm_{-n-1}(G,M)$), sedangkan untuk $n=0$ (atau $n=-1$), kohomologi
tersebut sama dengan $\Coker(\nu)$ (atau $\Ker(\nu)$). Kohomologi Tate tetap
memenuhi sifat-sifat seperti barisan eksak panjang dan Lema Shapiro. Dalam
banyak hal, teori ini bahkan lebih mudah dioperasikan daripada kohomologi atau
homologi dan mempunyai penerapan yang sangat menonjol dalam teori bilangan.

% segment-id: o014.li2.chapter6.overview-g-modules-resolutions.p014
Teori di atas hanya membahas grup diskret; kasus grup yang dilengkapi topologi
jauh lebih rumit. Grup profinit $G$ dan modul-$G$ mulus yang dibahas dalam
\S\ref{sec:profinite-cohomology} merupakan salah satu kasus khusus yang
sederhana sekaligus berguna; latar belakangnya terdapat dalam
\S\ref{sec:profinite-groups}. Pada dasarnya, teori ini tetap bersifat aljabar.
Kohomologi yang bersesuaian dapat dipahami sebagai funktor
turunan kanan, atau secara ekuivalen didefinisikan sebagai $\varinjlim$ dari
kohomologi grup hingga; versi kontinu kompleks standar menyediakan cara
perhitungan lain. Dalam teori bilangan, kasus utama yang diminati ialah ketika
$G$ merupakan grup Galois atau grup fundamental aritmetika.

% segment-id: o014.li2.chapter6.overview-g-modules-resolutions.p015
Terakhir, \S\ref{sec:nonabelian-coh} memperkenalkan kohomologi grup nonabelian.
Untuk grup nonabelian, biasanya hanya $\Hm^0$ dan $\Hm^1$ yang dapat
didefinisikan, sehingga barisan eksak panjangnya pun terbatas
(Teorema~\ref{prop:nonabelian-long}). Teori ini sangat sesuai untuk
mengklasifikasikan objek geometri atau aljabar, sebab grup simetri objek-objek
tersebut kerap nonabelian. Beberapa contoh akan dibahas kemudian dalam
\S\ref{sec:Galois-H1}. Konstruksi-konstruksi ini juga berlaku bagi grup
profinit.

% segment-id: o014.li2.chapter6.overview-g-modules-resolutions.n001
\begin{petunjukbacaan}
	Peran gelanggang koefisien $\Bbbk$ dalam bab ini bersifat sekunder dan sering
	diambil sama dengan $\Z$; lihat penjelasan dalam
	Catatan~\ref{rem:Gmod-k}. Pembaca dapat menyesuaikan dengan kebutuhannya dan
	melewati argumen yang melibatkan kategori turunan; langkah ini tidak akan
	mengganggu sebagian besar penerapan klasik. Secara khusus, hasil kali cawan
	pada kohomologi dapat diperlakukan sepenuhnya melalui operasi eksplisit
	dalam Definisi--Proposisi~\ref{def:group-cup-3}. Konsep modul-$G$ (lihat
	\S\ref{sec:G-mod}) sesekali muncul sebagai contoh atau latihan dalam bab-bab
	berikutnya dan biasanya hanya melibatkan definisi dasar. Pengecualiannya
	ialah \S\ref{sec:descent-Galois} dan \S\ref{sec:Galois-H1}, yang masing-masing
	memakai modul-$G$ mulus untuk grup profinit beserta kohomologinya. Selain itu,
	pembaca yang berminat pada teori bilangan sebaiknya menguasai isi
	\S\S\ref{sec:Tate-coh}---\ref{sec:nonabelian-coh}.
\end{petunjukbacaan}

% segment-id: o014.li2.chapter6.overview-g-modules-resolutions.h002
\section{Modul-\texorpdfstring{$G$}{G} dan Resolusinya}\label{sec:G-mod}

% segment-id: o014.li2.chapter6.overview-g-modules-resolutions.p016
Kecuali dinyatakan lain, dalam bagian ini $\Bbbk$ selalu merupakan gelanggang
komutatif dan $G$ suatu grup. Tuliskan unsur identitas grup $G$ sebagai $1_G$.

% segment-id: o014.li2.chapter6.overview-g-modules-resolutions.d001
\begin{definisi}\label{def:G-mod}
	\index{modul-$G$}
	\index[sym1]{G-Mod@$G\dcate{Mod}$}
	Suatu \emph{modul-$G$} dengan koefisien dalam $\Bbbk$ ialah modul-$\Bbbk$
	$M$ yang dilengkapi aksi kiri $G$. Aksi kiri $G\times M\to M$ ditulis
	sebagai perkalian dan memenuhi
% segment-id: o014.li2.chapter6.overview-g-modules-resolutions.q007
	\begin{align*}
		g(m + m') & = gm + gm', \\
		g(tm) & = t(gm),
	\end{align*}
	untuk sembarang $g\in G$, $t\in\Bbbk$, dan $m,m'\in M$. Dengan kata lain,
	$M$ dilengkapi aksi $G$ yang linear.

	Homomorfisme $f:M_1\to M_2$ antara dua modul-$G$ didefinisikan sebagai
	homomorfisme modul-$\Bbbk$ $f:M_1\to M_2$ yang memenuhi sifat ekuivarian
	$f(gm_1)=gf(m_1)$ untuk sembarang $g\in G$ dan $m_1\in M_1$.
\end{definisi}

% segment-id: o014.li2.chapter6.overview-g-modules-resolutions.p017
Secara tegas, struktur di atas seharusnya disebut modul-$G$ kiri. Modul-$G$
kanan dapat didefinisikan dengan cara yang sepenuhnya serupa dan sama artinya
dengan modul-$G^{\opp}$ kiri. Jika tidak menimbulkan kekeliruan, gelanggang
koefisien $\Bbbk$ dihilangkan dari notasi dan kategori semua modul-$G$ ditulis
$G\dcate{Mod}$. Kategori ini $\Bbbk$-linear
(Definisi~\ref{def:k-linear-cat}).

% segment-id: o014.li2.chapter6.overview-g-modules-resolutions.ex001
\begin{contoh}\label{eg:trivial-G-mod}
	\index{modul-$G$!trivial (trivial)}
	Lengkapi $\Bbbk$ dengan aksi trivial $G$, yaitu $gm=m$. Objek yang diperoleh
	disebut \emph{modul-$G$ trivial} dan cukup ditulis $\Bbbk$.
\end{contoh}

% segment-id: o014.li2.chapter6.overview-g-modules-resolutions.d002
\begin{definisi}
	Misalkan $M$ suatu modul-$G$. Jika submodul-$\Bbbk$ $N\subset M$ tertutup
	terhadap aksi $G$, maka $N$ disebut submodul-$G$. Lengkapi modul hasil bagi
	$M/N$ dengan aksi $G$ melalui $g(m+N)=gm+N$; objek ini disebut modul-$G$
	hasil bagi yang bersesuaian.
\end{definisi}

% segment-id: o014.li2.chapter6.overview-g-modules-resolutions.p018
Dari $G$ dapat dibangun aljabar grup $\Bbbk[G]$; lihat
\cite[Definisi 5.6.3]{Li1}. Setiap unsurnya dapat ditulis secara unik dalam
bentuk $\sum_{g\in G}a_gg$, dengan hanya berhingga banyak koefisien
$a_g\in\Bbbk$ yang tak nol. Mulai sekarang, $G$ dipandang sebagai himpunan
bagian dari $\Bbbk[G]$.

% segment-id: o014.li2.chapter6.overview-g-modules-resolutions.pr001
\begin{proposisi}\label{prop:G-mod-kG}
	Terdapat isomorfisme kategori
	$G\dcate{Mod}\rightiso\Bbbk[G]\dcate{Mod}$ sebagai berikut. Isomorfisme ini
	mempertahankan modul-$\Bbbk$ yang mendasarinya. Jika $M$ suatu modul-$G$,
	perkalian skalarnya sebagai modul kiri $\Bbbk[G]$ berasal dari homomorfisme
% segment-id: o014.li2.chapter6.overview-g-modules-resolutions.q008
	\[ \Bbbk[G] \dotimes{\Bbbk} M \to M, \quad (\sum_g a_g g) \otimes m \mapsto \sum_g a_g (gm), \]
	sedangkan homomorfisme modul-$G$ bersesuaian dengan homomorfisme modul
	$\Bbbk[G]$.
\end{proposisi}

% segment-id: o014.li2.chapter6.overview-g-modules-resolutions.pf001
\begin{bukti}
	Funktor invers ditentukan sebagai berikut. Untuk suatu modul kiri $\Bbbk[G]$
	$M$, gunakan pembenaman monoid multiplikatif $G\hookrightarrow\Bbbk[G]$
	untuk membuat $G$ beraksi dari kiri pada $M$, sehingga $M$ menjadi modul-$G$.
	Tanpa kesulitan dapat diverifikasi bahwa kedua arah terdefinisi dengan baik
	dan saling invers.
\end{bukti}

% segment-id: o014.li2.chapter6.overview-g-modules-resolutions.p019
Sebagai akibatnya, $G\dcate{Mod}$ merupakan kategori abelian $\Bbbk$-linear,
sebab demikian pula $\Bbbk[G]\dcate{Mod}$. Sifat-sifat submodul-$G$ dan
modul-$G$ hasil bagi cukup diterjemahkan ke kasus modul $\Bbbk[G]$.

% segment-id: o014.li2.chapter6.overview-g-modules-resolutions.d003
\begin{definisi}\label{def:HomG}
	\index[sym1]{HomG@$\Hom_G$}
	Tuliskan $\Hom$ dalam $G\dcate{Mod}$ sebagai $\Hom_G$. Di bawah korespondensi
dalam Proposisi~\ref{prop:G-mod-kG}, objek ini juga sama dengan
	$\Hom_{\Bbbk[G]}$ dalam $\Bbbk[G]\dcate{Mod}$.
\end{definisi}

% segment-id: o014.li2.chapter6.overview-g-modules-resolutions.p020
Beberapa operasi dasar pada modul-$\Bbbk$ mudah diangkat ke konteks modul-$G$:
% segment-id: o014.li2.chapter6.overview-g-modules-resolutions.l001
\begin{description}
% segment-id: o014.li2.chapter6.overview-g-modules-resolutions.i001
	\item[Jumlah langsung dan produk langsung] Untuk suatu keluarga modul-$G$
	$(M_i)_{i\in I}$, lengkapi modul-$\Bbbk$ $\prod_{i\in I}M_i$ dan
	$\bigoplus_{i\in I}M_i$ dengan aksi kiri $G$
% segment-id: o014.li2.chapter6.overview-g-modules-resolutions.q009
	\[ g (m_i)_{i \in I} := (gm_i)_{i \in I}, \]
	yang disebut aksi diagonal. Dengan aksi ini, $\prod_{i\in I}M_i$ dan
	$\bigoplus_{i\in I}M_i$ menjadi modul-$G$.
% segment-id: o014.li2.chapter6.overview-g-modules-resolutions.i002
	\item[Hasil kali tensor] Untuk modul-$G$ $M$ dan $N$, lengkapi
	$M\otimes N:=M\dotimes{\Bbbk}N$ dengan aksi kiri diagonal $G$:
% segment-id: o014.li2.chapter6.overview-g-modules-resolutions.q010
	\[ g (x \otimes y) = gx \otimes gy, \quad x \in M, \quad y \in N ; \]
	karena $(x,y)\mapsto gx\otimes gy$ merupakan pemetaan bilinear, rumus ini
	terdefinisi dengan baik. Jelas bahwa
	$x\otimes y\mapsto y\otimes x$ mendefinisikan isomorfisme modul-$G$
	$M\otimes N\simeq N\otimes M$.
% segment-id: o014.li2.chapter6.overview-g-modules-resolutions.i003
	\item[Modul $\Hom$] Untuk modul-$G$ $M$ dan $N$, homomorfisme
	modul-$\Bbbk$ di antaranya membentuk modul-$\Bbbk$
	$\Hom(M,N):=\Hom_{\Bbbk}(M,N)$. Lengkapi $\Hom(M,N)$ dengan aksi kiri $G$
% segment-id: o014.li2.chapter6.overview-g-modules-resolutions.q011
	\[ ({}^g f)(x) = g(f(g^{-1} x)), \quad f \in \Hom(M, N), \quad x \in M, \]
	dengan $g^{-1}$ dan $g$ di ruas kanan masing-masing berasal dari aksi kiri
	$G$ pada $M$ dan $N$. Sebagai komposisi tiga homomorfisme modul-$\Bbbk$,
	${}^gf$ berada dalam $\Hom(M,N)$. Dengan demikian, ${}^gf$ dapat dipahami
	sebagai ``konjugasi'' $gfg^{-1}$, dan sifat-sifat yang diperlukan untuk
	struktur modul-$G$ pun menjadi jelas.

% segment-id: o014.li2.chapter6.overview-g-modules-resolutions.i004
	\item[Modul kontragradien] Dalam konstruksi modul $\Hom$, ambil
	$N=\Bbbk$ sebagai modul-$G$ trivial. Dengan demikian, $\Hom(M,\Bbbk)$
	menjadi modul-$G$ yang disebut modul kontragradien dari $M$. Aksi $G$ padanya
	diberikan oleh ${}^gf=f\circ g^{-1}$.
	\index{modul-$G$!kontragradien (contragredient)}
\end{description}

% segment-id: o014.li2.chapter6.overview-g-modules-resolutions.p021
Semua konstruksi ini funktorial terhadap $M$ dan $N$. Definisi tersebut segera
memberikan
% segment-id: o014.li2.chapter6.overview-g-modules-resolutions.q012
\begin{align*}
	\Hom_G(M, N) & = \Hom(M, N)^G \\
	& := \left\{f \in \Hom(M, N): \forall g \in G, \; {}^g f = f \right\}.
\end{align*}

% segment-id: o014.li2.chapter6.overview-g-modules-resolutions.p022
Sehubungan dengan isomorfisme kategori dalam
Proposisi~\ref{prop:G-mod-kG}, jumlah langsung dan produk langsung modul-$G$
masing-masing ialah jumlah langsung dan produk langsung modul $\Bbbk[G]$.
Struktur modul-$G$ pada hasil kali tensor dan modul $\Hom$ sedikit lebih rumit;
dari sudut pandang teori modul, struktur tersebut sesungguhnya bersesuaian
dengan struktur aljabar Hopf pada $\Bbbk[G]$
(Contoh~\ref{eg:group-Hopf-algebra}). Pembaca yang berminat dapat
membandingkannya dengan pembahasan terkait dalam \S\ref{sec:bialgebra} dan
\S\ref{sec:duality-examples}; kita tidak akan menguraikannya di sini.

% segment-id: o014.li2.chapter6.overview-g-modules-resolutions.pr002
\begin{proposisi}\label{prop:Gmod-monoidal}
	Terhadap operasi hasil kali tensor, $G\dcate{Mod}$ menjadi kategori monoidal
	simetris dengan modul-$G$ trivial $\Bbbk$ sebagai objek unit.
\end{proposisi}

% segment-id: o014.li2.chapter6.overview-g-modules-resolutions.pf002
\begin{bukti}
	Ini merupakan akibat langsung dari definisi di atas.
\end{bukti}

% segment-id: o014.li2.chapter6.overview-g-modules-resolutions.p023
Selain itu, isomorfisme kanonik yang lazim dalam $\Bbbk\dcate{Mod}$,
% segment-id: o014.li2.chapter6.overview-g-modules-resolutions.g001
\begin{equation*}
	\begin{tikzcd}[row sep=tiny]
		\Hom(L \otimes M, N) \arrow[leftrightarrow, r, "\sim"] & \Hom(L, \Hom(M, N)) \\
		\varphi \arrow[mapsto, r] & {\left[ w \mapsto \varphi(w \otimes (\cdot)) \right]} \\
		{\left[ w \otimes x \mapsto \psi(w)(x) \right]} & \psi \arrow[mapsto, l]
	\end{tikzcd}
\end{equation*}
juga terangkat ke tingkat modul-$G$. Jika $L$, $M$, dan $N$ merupakan
modul-$G$, dan hasil kali tensor serta $\Hom$ di atas dilengkapi struktur
modul-$G$ yang telah didefinisikan, maka kedua arah merupakan isomorfisme
modul-$G$; pembaca dipersilakan memverifikasinya dengan teliti.

% segment-id: o014.li2.chapter6.overview-g-modules-resolutions.p024
Sebagai akibatnya, ${}^g\varphi=\varphi$ ekuivalen dengan
${}^g\psi=\psi$, sehingga diperoleh adjungsi
% segment-id: o014.li2.chapter6.overview-g-modules-resolutions.e001
\begin{equation}\label{eqn:G-mod-tensor-Hom}
	\begin{tikzcd}
		\Hom_G(L \otimes M, N) \arrow[leftrightarrow, r, "\sim"] & \Hom_G(L, \Hom(M, N)),
	\end{tikzcd}
\end{equation}
dengan $L$, $M$, dan $N$ semuanya modul-$G$. Berikut sebuah penerapan
sederhana.

% segment-id: o014.li2.chapter6.overview-g-modules-resolutions.pr003
\begin{proposisi}\label{prop:group-tensor-proj}
	Misalkan $L$ dan $M$ modul-$G$, dengan $L$ projektif sebagai modul-$G$ dan
	$M$ projektif sebagai modul-$\Bbbk$. Maka $L\otimes M$ projektif sebagai
	modul-$G$.
\end{proposisi}

% segment-id: o014.li2.chapter6.overview-g-modules-resolutions.pf003
\begin{bukti}
	Adjungsi \eqref{eqn:G-mod-tensor-Hom} menunjukkan bahwa funktor
	$(\cdot)\otimes M:G\dcate{Mod}\to G\dcate{Mod}$ mempunyai adjoin kanan eksak
	$\Hom(M,\cdot)$. Karena itu,
	Proposisi~\ref{prop:adjoint-injective-projective} menyiratkan bahwa
	$(\cdot)\otimes M$ mempertahankan objek projektif.
\end{bukti}

% segment-id: o014.li2.chapter6.overview-g-modules-resolutions.d004
\begin{definisi}[Restriksi dan inflasi]\label{def:Res-Infl}
	\index{modul-$G$!restriksi/inflasi (restriction/inflation)}
	\index[sym1]{ResGH@$\Res^G_H$}
	\index[sym1]{InflGH@$\mathrm{Infl}^G_{G/H}$}
	Untuk setiap homomorfisme grup $\varphi:H\to G$, terdapat funktor yang
	bersesuaian
% segment-id: o014.li2.chapter6.overview-g-modules-resolutions.q013
	\[ \varphi^*: G\dcate{Mod} \to H\dcate{Mod}; \]
	bagi modul-$G$ $M$, citranya oleh $\varphi^*$ ialah modul-$\Bbbk$ yang sama,
	$M$, dengan aksi $H$ yang diberikan oleh $hm:=\varphi(h)m$. Dua kasus penting
	dari funktor $\varphi^*$ ialah:
% segment-id: o014.li2.chapter6.overview-g-modules-resolutions.l002
	\begin{description}
% segment-id: o014.li2.chapter6.overview-g-modules-resolutions.i005
		\item[Restriksi] Untuk subgrup $H\subset G$, ambil
		$\varphi:H\hookrightarrow G$. Funktor yang bersesuaian
		$\Res^G_H:G\dcate{Mod}\to H\dcate{Mod}$ membatasi aksi grup pada $H$.
% segment-id: o014.li2.chapter6.overview-g-modules-resolutions.i006
		\item[Inflasi] Untuk subgrup normal $H\lhd G$, ambil
		$\varphi:G\twoheadrightarrow G/H$. Funktor yang bersesuaian
		$\mathrm{Infl}^G_{G/H}:(G/H)\dcate{Mod}\to G\dcate{Mod}$ menarik balik
		aksi grup ke $G$.
	\end{description}
\end{definisi}

% segment-id: o014.li2.chapter6.overview-g-modules-resolutions.p025
Perhatikan bahwa jika terdapat rangkaian homomorfisme grup
$H\xrightarrow{\varphi}G\xrightarrow{\psi}F$, maka funktor-funktor yang
bersesuaian memenuhi $\varphi^*\psi^*=(\psi\varphi)^*$.

% segment-id: o014.li2.chapter6.overview-g-modules-resolutions.p026
Jika sudut pandang aljabar grup dipakai, homomorfisme $\varphi:H\to G$ meluas
menjadi homomorfisme aljabar-$\Bbbk$ $\Bbbk[H]\to\Bbbk[G]$, dan funktor yang
bersesuaian $\Bbbk[G]\dcate{Mod}\to\Bbbk[H]\dcate{Mod}$ tepat bersesuaian
dengan $\varphi^*$. Funktor ini juga dapat dinyatakan melalui hasil kali tensor
sebagai
% segment-id: o014.li2.chapter6.overview-g-modules-resolutions.e002
\begin{equation}\label{eqn:G-module-pullback-tensor}
	\varphi^* M \simeq \Bbbk[G] \dotimes{\Bbbk[G]} M, \quad \Bbbk[G] \;\text{melalui $\varphi$ dipandang sebagai bimodul}\; (\Bbbk[H], \Bbbk[G]).
\end{equation}

% segment-id: o014.li2.chapter6.overview-g-modules-resolutions.p027
Adjoin kiri dan adjoin kanan dari funktor $\varphi^*$, beserta isomorfisme
adjungsi yang bersesuaian, dapat dituliskan secara eksplisit.

% segment-id: o014.li2.chapter6.overview-g-modules-resolutions.pr004
\begin{proposisi}\label{prop:pullback-two-adjoint}
	Gunakan notasi di atas. Untuk suatu homomorfisme grup $\varphi:H\to G$,
	sembarang modul-$G$ $M$, dan sembarang modul-$H$ $N$, terdapat isomorfisme
	kanonik modul-$\Bbbk$ berikut:
% segment-id: o014.li2.chapter6.overview-g-modules-resolutions.g002
	\[\begin{tikzcd}[row sep=tiny]
		\Hom_H(\varphi^* M, N) \arrow[leftrightarrow, r, "\sim"] & \Hom_G\left( M, \Hom_H(\Bbbk[G], N) \right) \\
		\theta \arrow[mapsto, r] & {[x \mapsto [G \ni g \mapsto \theta(gx)] ]} \\
		{[x \mapsto \xi(x)(1_G)]} & \xi, \arrow[mapsto, l] \\
		\Hom_H(N, \varphi^* M) \arrow[leftrightarrow, r, "\sim"] & \Hom_G\left( \Bbbk[G] \dotimes{\Bbbk[H]} N, M \right) \\
		\psi \arrow[mapsto, r] & {[g \otimes y \mapsto g\psi(y)]} \\
		{[y \mapsto \eta(1_G \otimes y) ]} & \eta \arrow[mapsto, l].
	 \end{tikzcd}\]
	Di sini, $\Bbbk[G]$ dipandang sekaligus sebagai bimodul
	$(\Bbbk[H],\Bbbk[G])$ dan bimodul $(\Bbbk[G],\Bbbk[H])$. Dengan cara ini,
	$\Hom_H(\Bbbk[G],N)$ dan $\Bbbk[G]\dotimes{\Bbbk[H]}N$ dilengkapi struktur
	modul-$G$.
\end{proposisi}

% segment-id: o014.li2.chapter6.overview-g-modules-resolutions.pf004
\begin{bukti}
	Uraikan langsung definisinya untuk memverifikasi bahwa kedua pasang pemetaan
	merupakan homomorfisme modul-$\Bbbk$ yang terdefinisi dengan baik;
	perhitungannya tidak rumit. Setelah hal ini diverifikasi, sifat saling invers
	menjadi jelas. Cara lain ialah menerapkan teori umum dalam
	\cite[Korolari 6.6.8]{Li1}.
\end{bukti}

% segment-id: o014.li2.chapter6.overview-g-modules-resolutions.p028
Karena $G\dcate{Mod}$ kategori abelian, resolusi projektif dan resolusi
injektif modul-$G$ dapat dibahas secara alami. Untuk tema bab ini, yang paling
penting ialah resolusi projektif dari modul-$G$ trivial $\Bbbk$. Resolusi
tersebut tidak hanya ada, tetapi juga mempunyai pilihan yang kanonik dan
eksplisit.

% segment-id: o014.li2.chapter6.overview-g-modules-resolutions.d005
\begin{definisi}\label{def:resolution-trivial-module}
	Misalkan $G$ suatu grup. Untuk setiap $n\in\Z_{\geq0}$, definisikan
% segment-id: o014.li2.chapter6.overview-g-modules-resolutions.q014
	\[ \mathsf{L}_n := \text{modul-$\Bbbk$ bebas dengan basis $G^{n+1}$}, \]
	yang unsur-unsurnya ditulis sebagai kombinasi linear-$\Bbbk$ dari
	$(g_0,\ldots,g_n)\in G^{n+1}$. Untuk setiap $n\geq1$, definisikan
	homomorfisme
% segment-id: o014.li2.chapter6.overview-g-modules-resolutions.q015
	\begin{align*}
		\partial_n, \partial'_n: \mathsf{L}_n & \to \mathsf{L}_{n-1}, \\
		\partial_n\left( g_0, \ldots, g_n \right) & = \sum_{k=0}^n (-1)^{n-k} (\ldots, \widehat{g_k}, \ldots), \\
		\partial'_n\left( g_0, \ldots, g_n \right) & = \sum_{k=0}^n (-1)^k (\ldots, \widehat{g_k}, \ldots),
	\end{align*}
	dengan $\widehat{g_k}$ menandakan bahwa suku tersebut dihilangkan; jadi,
	$\partial'_n=(-1)^n\partial_n$. Definisikan pula homomorfisme
% segment-id: o014.li2.chapter6.overview-g-modules-resolutions.q016
	\[ \partial_0 = \partial'_0: \mathsf{L}_0 \to \Bbbk, \quad (g_0) \mapsto 1. \]

	Definisikan aksi
	$g(g_0,\ldots,g_n)=(gg_0,\ldots,gg_n)$ (atau
	$(g_0,\ldots,g_n)g=(g_0g,\ldots,g_ng)$) agar $\mathsf{L}_n$ menjadi
	modul-$G$ kiri (atau kanan), dan lengkapi $\Bbbk$ dengan aksi trivial $G$.
	Dengan demikian, setiap $\partial_n$ dan $\partial'_n$ merupakan
	homomorfisme modul-$G$ kiri (atau kanan).
\end{definisi}

% segment-id: o014.li2.chapter6.overview-g-modules-resolutions.p029
Jelas bahwa $(\mathsf{L}_n,\partial_n)_n$ dan
$(\mathsf{L}_n,\partial'_n)_n$ hanya berbeda melalui pembalikan
$(g_0,\ldots,g_n)\mapsto(g_n,\ldots,g_0)$. Pembalikan ini berkomutasi dengan
aksi kiri maupun kanan $G$, sehingga cukup ditangani salah satu dari
$\partial_n$ dan $\partial'_n$ untuk semua sifat berikut.

% segment-id: o014.li2.chapter6.overview-g-modules-resolutions.r001
\begin{catatan}\label{rem:augmentation-kG}
	\index{homomorfisme augmentasi (augmentation homomorphism)}
	Jika $\mathsf{L}_0$ diidentifikasi dengan $\Bbbk[G]$, maka
	$\partial_0=\partial'_0:\mathsf{L}_0\to\Bbbk$ diidentifikasi dengan
	pemetaan
% segment-id: o014.li2.chapter6.overview-g-modules-resolutions.q017
	\[ \Bbbk[G] \to \Bbbk, \quad \sum_g a_g g \mapsto \sum_g a_g; \]
	jelas bahwa pemetaan ini juga merupakan homomorfisme surjektif aljabar
	$\Bbbk$. Pemetaan tersebut disebut \emph{homomorfisme augmentasi} dari
	$\Bbbk[G]$, dan kernelnya
% segment-id: o014.li2.chapter6.overview-g-modules-resolutions.q018
	\[ \mathfrak{I} := \left\{ \sum_g a_g g : \sum_g a_g = 0 \right\} \]
	disebut \emph{ideal augmentasi} dari $\Bbbk[G]$. Sebagai modul-$\Bbbk$,
	$\mathfrak{I}$ jelas terurai terhadap basis berikut:
% segment-id: o014.li2.chapter6.overview-g-modules-resolutions.q019
	\[ \mathfrak{I} = \bigoplus_{\substack{g \in G \\ g \neq 1_G}} \Bbbk (g - 1_G). \]
	\index{ideal augmentasi (augmentation ideal)}
	\index[sym1]{I-aug@$\mathfrak{I}$}
\end{catatan}

% segment-id: o014.li2.chapter6.overview-g-modules-resolutions.le001
\begin{lema}\label{prop:L-acyclic}
	Kompleks
	$\cdots\to\mathsf{L}_n\xrightarrow{\partial_n}\cdots
	\xrightarrow{\partial_0}\Bbbk\to0$ yang didefinisikan di atas eksak. Jika
	dipandang sebagai kompleks modul-$\Bbbk$, morfisme identitasnya homotopik nol.
	Pernyataan yang sama berlaku jika $\partial_n$ diganti dengan $\partial'_n$.
\end{lema}

% segment-id: o014.li2.chapter6.overview-g-modules-resolutions.pf005
\begin{bukti}
	Cukup bahas versi $\partial'_n$. Untuk setiap $n\geq1$, perhatikan bahwa
	$\partial'_{n-1}\partial'_n$ memetakan $(g_0,\ldots,g_n)$ ke suatu jumlah
	unsur berbentuk
% segment-id: o014.li2.chapter6.overview-g-modules-resolutions.q020
	\[ \pm (\ldots, \widehat{g_p}, \ldots, \widehat{g_q}, \ldots), \quad 0 \leq p < q \leq n. \]
	Tanda-tandanya saling meniadakan berpasangan, sehingga diperoleh suatu
	kompleks; verifikasi sederhananya diserahkan kepada pembaca.

	Selanjutnya, untuk setiap $n\geq0$, definisikan homomorfisme modul-$\Bbbk$
% segment-id: o014.li2.chapter6.overview-g-modules-resolutions.q021
	\[ s_n: \mathsf{L}_n \to \mathsf{L}_{n+1}, \quad (g_0, \ldots, g_n) \mapsto (1_G, g_0, \ldots, g_n), \]
	dan definisikan $s_{-1}:\Bbbk\to\mathsf{L}_0$ dengan memetakan
	$1\in\Bbbk$ ke $1_G\in\mathsf{L}_0$. Mudah dilihat bahwa
% segment-id: o014.li2.chapter6.overview-g-modules-resolutions.q022
	\[ \partial'_{n+1} s_n + s_{n-1} \partial'_n = \identity_{\mathsf{L}_n} \]
	untuk setiap $n\geq-1$, dengan konvensi $\mathsf{L}_{-1}:=\Bbbk$ dan
	$\partial'_{-1}:=0$. Hal ini menunjukkan bahwa
% segment-id: o014.li2.chapter6.overview-g-modules-resolutions.q023
	\[ \cdots \to \mathsf{L}_n \xrightarrow{\partial'_n} \cdots \xrightarrow{\partial'_0} \Bbbk \to 0 \]
	memiliki morfisme identitas yang homotopik nol sebagai kompleks
	modul-$\Bbbk$, sehingga kompleks tersebut eksak.
\end{bukti}

% segment-id: o014.li2.chapter6.overview-g-modules-resolutions.x001
\index[sym1]{L-cplx@$\mathsf{L}$, $\overline{\mathsf{L}}$}

% segment-id: o014.li2.chapter6.overview-g-modules-resolutions.p030
Tuliskan kompleks
$\cdots\to\mathsf{L}_n\xrightarrow{\partial_n}\cdots\to\mathsf{L}_0\to0$
sebagai $\mathsf{L}$. Kompleks ini dapat dipandang sebagai kompleks pada
derajat $\ldots,-n,\ldots,0$, atau sebagai kompleks rantai pada derajat
$\ldots,n,\ldots,0$. Dalam kedua pandangan tersebut,
$\partial_0=\partial'_0$ memberikan morfisme $\mathsf{L}\to\Bbbk$.

% segment-id: o014.li2.chapter6.overview-g-modules-resolutions.p031
Kompleks dan morfisme ini dapat dipahami baik sebagai kompleks modul-$G$ kiri
maupun sebagai kompleks modul-$G$ kanan. Berdasarkan simetri, kita terutama akan
menyatakan versi kiri.

% segment-id: o014.li2.chapter6.overview-g-modules-resolutions.pr005
\begin{proposisi}\label{prop:trivial-mod-resolution}
	Morfisme $\mathsf{L}\to\Bbbk$ yang didefinisikan di atas memberikan resolusi
	bebas dari modul-$G$ trivial $\Bbbk$.
\end{proposisi}

% segment-id: o014.li2.chapter6.overview-g-modules-resolutions.pf006
\begin{bukti}
	Ketika $(h_1,\ldots,h_n)\in G^n$ bervariasi,
	$(1_G,h_1,\ldots,h_n)$ merentang suatu basis dari $\mathsf{L}^n$ sebagai
	modul kiri $\Bbbk[G]$.
\end{bukti}

% segment-id: o014.li2.chapter6.overview-g-modules-resolutions.p032
Ada pula resolusi bebas kanonik dari $\Bbbk$ yang lebih ringkas dan kadang-kadang
lebih praktis daripada $\mathsf{L}$.

% segment-id: o014.li2.chapter6.overview-g-modules-resolutions.dp001
\begin{definisi-proposisi}[Kompleks ternormalisasi]
	Untuk setiap $n>0$, definisikan $\mathrm{Triv}_n\subset\mathsf{L}_n$ sebagai
	submodul-$G$ yang dihasilkan oleh unsur-unsur
% segment-id: o014.li2.chapter6.overview-g-modules-resolutions.q024
	\[ (g_0, \ldots, g_n) \in G^{n+1}, \quad \exists\, 0 \leq k < n, \; g_k = g_{k+1}. \]
	Definisikan pula $\mathrm{Triv}_0=0$. Baik terhadap $\partial_n$ maupun
	$\partial'_n$, objek-objek ini membentuk subkompleks $\mathrm{Triv}$ dari
	$\mathsf{L}$. Normalisasi $\mathsf{L}$ kemudian didefinisikan sebagai
	kompleks hasil bagi
	$\overline{\mathsf{L}}:=\mathsf{L}/\mathrm{Triv}$. Jika dipandang sebagai
	kompleks modul-$\Bbbk$, morfisme identitas pada $\mathrm{Triv}$ homotopik nol.
\end{definisi-proposisi}

% segment-id: o014.li2.chapter6.overview-g-modules-resolutions.pf007
\begin{bukti}
	Jelas bahwa $\mathrm{Triv}_n$ invarian terhadap
	$(g_0,\ldots,g_n)\mapsto(g_n,\ldots,g_0)$, sehingga versi $\partial_n$ dan
	$\partial'_n$ tidak berbeda; cukup ditangani versi terakhir. Jika
	$(g_0,\ldots,g_n)$ memenuhi $g_k=g_{k+1}$, maka dalam
	$\partial'_n(g_0,\ldots,g_n)=
	\sum_i(-1)^i(\ldots,\widehat{g_i},\ldots)$, suku-suku yang bersesuaian dengan
	$i=k,k+1$ saling meniadakan, sedangkan suku-suku lain berada dalam
	$\mathrm{Triv}_{n-1}$. Ini membuktikan bahwa objek-objek tersebut membentuk
	subkompleks.

	Untuk menunjukkan bahwa morfisme identitas $\mathrm{Triv}$ homotopik nol
	sebagai kompleks modul-$\Bbbk$, cukup batasi homotopi
	$s_n:\mathsf{L}_n\to\mathsf{L}_{n+1}$ dalam pembuktian
	Lema~\ref{prop:L-acyclic} pada $\mathrm{Triv}_n$, untuk $n\geq0$.
\end{bukti}

% segment-id: o014.li2.chapter6.overview-g-modules-resolutions.p033
Berdasarkan konstruksi, $\mathsf{L}\to\Bbbk$ bernilai nol pada
$\mathrm{Triv}$, sehingga memfaktor melalui
$\overline{\mathsf{L}}\to\Bbbk$.

% segment-id: o014.li2.chapter6.overview-g-modules-resolutions.pr006
\begin{proposisi}\label{prop:trivial-mod-nor-resolution}
	Morfisme $\overline{\mathsf{L}}\to\Bbbk$ yang didefinisikan di atas
	memberikan resolusi bebas dari modul-$G$ trivial $\Bbbk$.
\end{proposisi}

% segment-id: o014.li2.chapter6.overview-g-modules-resolutions.pf008
\begin{bukti}
	Pertama, setiap
	$\overline{\mathsf{L}}_n=\mathsf{L}_n/\mathrm{Triv}_n$ merupakan modul-$G$
	bebas: suatu basis diberikan oleh citra
	$(1_G,h_1,\ldots,h_n)\in G^{n+1}$ dengan $h_1\neq1_G$ dan
	$h_k\neq h_{k+1}$ untuk $1\leq k<n$.

	Kedua, keasiklikan $\mathrm{Triv}$ menyiratkan bahwa
	$\mathsf{L}\to\overline{\mathsf{L}}$ merupakan kuasi-isomorfisme. Karena
	itu, Proposisi~\ref{prop:trivial-mod-resolution} menyiratkan bahwa
	$\overline{\mathsf{L}}\to\Bbbk$ merupakan kuasi-isomorfisme.
\end{bukti}

% segment-id: o014.li2.chapter6.overview-g-modules-resolutions.p034
Untuk penerapan selanjutnya, akan lebih praktis menyatakan basis dalam bentuk
``bar'' berikut. Pertama, untuk bilangan bulat $a\leq b$ dan suatu rangkaian
unsur $(g_a,\ldots,g_b)\in G^{b-a+1}$, perkenalkan notasi singkat
% segment-id: o014.li2.chapter6.overview-g-modules-resolutions.q025
\[ g_{[a, b]} := g_a g_{a+1} \cdots g_b. \]
Sebagai modul-$G$ kiri, basis $\mathsf{L}_n$ dapat diambil sebagai
% segment-id: o014.li2.chapter6.overview-g-modules-resolutions.e003
\begin{equation}\label{eqn:L-basis}
	(g_1 | \cdots | g_n) := \left( 1_G, g_{[1, 1]}, g_{[1, 2]}, \ldots, g_{[1, n]} \right), \quad (g_1, \ldots, g_n) \in G^n.
\end{equation}

% segment-id: o014.li2.chapter6.overview-g-modules-resolutions.p035
Perhatikan bahwa
% segment-id: o014.li2.chapter6.overview-g-modules-resolutions.e004
\begin{equation}\label{eqn:L-basis-partial}\begin{aligned}
	\partial_n (g_1 | \ldots | g_n) & = (-1)^n \left( g_{[1, 1]}, \ldots, g_{[1, n]}\right) + (-1)^{n-1} \left(1_G, g_{[1, 2]}, \ldots, g_{[1, n]}\right) \\
	& \quad + \cdots + \left( 1_G, g_{[1, 1]}, \ldots, g_{[1, n-1]}\right) \\
	& = (-1)^n g_1 (g_2 | \cdots | g_n) + \sum_{k=1}^{n-1} (-1)^{n-k} (\cdots | g_k g_{k+1}| \cdots) \\
	& \quad + (g_1 | \cdots |g_{n-1}), \\
	\partial'_n(g_1 | \cdots | g_n) & = g_1 (g_2 | \cdots | g_n) + \sum_{k=1}^{n-1} (-1)^k (\cdots | g_k g_{k+1}| \cdots) \\
	& \quad + (-1)^n (g_1 | \cdots |g_{n-1}).
\end{aligned}\end{equation}

% segment-id: o014.li2.chapter6.overview-g-modules-resolutions.p036
Sekarang tinjau bayangan cermin dari \eqref{eqn:L-basis}. Ambil basis
$\mathsf{L}_n$ sebagai modul-$G$ kanan berupa
% segment-id: o014.li2.chapter6.overview-g-modules-resolutions.e005
\begin{equation}\label{eqn:L-basis-right}
	(g_1 | \cdots | g_n)^\dagger := \left( g_{[1, n]}, \ldots, g_{[n-1, n]} , g_{[n, n]}, 1_G \right), \quad (g_1, \ldots, g_n) \in G^n.
\end{equation}
Dengan demikian,
$\partial'_n(g_1|\cdots|g_n)^\dagger=
(-1)^n\partial_n(g_1|\cdots|g_n)^\dagger$ sama dengan
% segment-id: o014.li2.chapter6.overview-g-modules-resolutions.e006
\begin{multline}\label{eqn:L-basis-partial-right}
	\left( g_{[2, n]}, \ldots, g_{[n, n]}, 1_G \right) - \left( g_{[1, n]}, g_{[3, n]}, \ldots, g_{[n, n]}, 1_G \right) + \cdots \\
	+ (-1)^n \left( g_{[1, n-1]}, \ldots, g_{[n-1, n-1]} \right) g_n  \\
	= (g_2 | \cdots | g_n)^\dagger + \sum_{k=1}^{n-1} (-1)^k (\cdots | g_k g_{k+1} | \cdots)^\dagger \\
	+ (-1)^n (g_1 | \cdots | g_{n-1})^\dagger g_n.
\end{multline}

% segment-id: o014.li2.chapter6.overview-g-modules-resolutions.p037
Kasus kompleks ternormalisasi mengikuti dengan cara yang sama: basis
$\overline{\mathsf{L}}_n$ sebagai modul-$G$ kiri atau kanan masing-masing
dapat diambil berupa
% segment-id: o014.li2.chapter6.overview-g-modules-resolutions.e007
\begin{equation}\label{eqn:nor-cochain}
	(g_1| \cdots | g_n) \quad \text{atau} \quad (g_1 | \cdots | g_n)^\dagger \quad \bmod \mathrm{Triv}_n,
\end{equation}
dengan syarat $g_k\neq1_G$ untuk setiap $1\leq k\leq n$.

% segment-id: o014.li2.chapter6.overview-g-modules-resolutions.r002
\begin{catatan}
	Resolusi di atas mempunyai penafsiran topologis yang mendalam:
	$\mathsf{L}$ ialah kompleks rantai yang diberikan oleh torsor universal
	$\mathrm{E}G$ di atas ruang pengklasifikasi $\mathrm{B}G$; lihat
	Contoh~\ref{eg:classifying-space-std-cplx} dan
	\ref{eg:classifying-contractible}.
\end{catatan}
